\emph{Preliminaries.} The blades are a basis of $\Cl_{\Q}(\ell,\ell)$ (Theorem \claimref{synthesis.signed-directions-and-the-integer-clifford-algebra}), so its dimension is $4^{\ell}$; the adjoint $\dagger$ is the linear extension of the inverse on blades, an anti-automorphism with $e_{j}^{\dagger} = e_{j}$, $f_{j}^{\dagger} = -f_{j}$, hence $v_{j}^{\dagger} = w_{j}$. By Proposition \claimref{tower.ladder-pairs-satisfy-car-vectors} (the relations are also checked directly for three modes), $Q(v_{j}) = Q(w_{j}) = 0$, $v_{j}w_{j} + w_{j}v_{j} = 1$, and ladder elements of different modes anticommute.

\emph{(i).} Put $p_{j} = w_{j}v_{j}$. Then $p_{j}^{2} = w_{j}(1 - w_{j}v_{j})v_{j} = p_{j}$, $p_{j}^{\dagger} = p_{j}$, the $p_{j}$ are even and commute with each other and with the ladder elements of other modes, $w_{j}p_{j} = 0 = p_{j}v_{j}$, and $p_{j} = \tfrac12(1 + e_{j}f_{j})$. So $P_{0} = \prod_{j}p_{j}$ is a self-adjoint idempotent with $\langle P_{0}\rangle_{0} = 2^{-\ell} \ne 0$, $w_{j}P_{0} = 0$ and $P_{0}v_{j} = 0$ for every $j$. For subsets $S, T$: $s_{S}^{\dagger}s_{T} = P_{0}w_{j_{k}}\cdots w_{j_{1}}v_{t_{1}}\cdots v_{t_{n}}P_{0}$. If some $t \in T$ is not in $S$, move $v_{t}$ to the left (it anticommutes with every other factor) until it meets $P_{0}$: the product is $0$. If some $j \in S$ is not in $T$, move $w_{j}$ to the right: again $0$. If $S = T$, the innermost pair is $w_{j_{1}}v_{j_{1}} = 1 - v_{j_{1}}w_{j_{1}}$, and in the second term $w_{j_{1}}$ moves to the right up to $P_{0}$ and gives $0$; by induction on $k$ the product is $P_{0}$. So $s_{S}^{\dagger}s_{T} = \delta_{ST}P_{0}$. The elements $E_{ST} = s_{S}s_{T}^{\dagger}$ then satisfy $E_{ST}E_{UV} = \delta_{TU}E_{SV}$ and $s_{S}^{\dagger}E_{ST}s_{T} = P_{0} \ne 0$. They are linearly independent (multiply a relation $\sum c_{ST}E_{ST} = 0$ by $E_{UU}$ on the left and $E_{VV}$ on the right) and there are $4^{\ell}$ of them, so they are a basis. Writing $1 = \sum c_{ST}E_{ST}$ and multiplying by $E_{UV}$ on the right gives $c_{SU} = \delta_{SU}$, so $\sum_{S}s_{S}s_{S}^{\dagger} = 1$. As in the theorem for two modes, $m(x)_{ST}$ defined by $s_{S}^{\dagger}xs_{T} = m(x)_{ST}P_{0}$ (well defined, because $s_{S}^{\dagger}xs_{T} = \sum c_{UV}s_{S}^{\dagger}E_{UV}s_{T} = c_{ST}P_{0}$ for $x = \sum c_{UV}E_{UV}$) is multiplicative, since $s_{S}^{\dagger}xys_{T} = \sum_{U}(s_{S}^{\dagger}xs_{U})(s_{U}^{\dagger}ys_{T})$, bijective, and $s_{S}^{\dagger}x^{\dagger}s_{T} = (s_{T}^{\dagger}xs_{S})^{\dagger} = m(x)_{TS}P_{0}$. The element $s_{S}$ has the parity of $|S|$, so for even $x$ the element $s_{S}^{\dagger}xs_{T}$ has the parity of $|S| + |T|$; it is a multiple of the even element $P_{0}$, so it vanishes when $|S| + |T|$ is odd. The matrices that vanish between subsets of different parity have dimension $2 \cdot 4^{\ell-1}$, the dimension of the even subalgebra, so the image of the even subalgebra is all of them. The check verifies every statement of this clause for $\ell = 3$ on all $64$ blades.

\emph{(ii).} By \claimref{std-algebra.quaternions} the even subalgebra of $\Cl_{\Q}(3,0)$ is spanned by $1, B_{1}, B_{2}, B_{3}$ with $B_{k}^{2} = -1$ and $B_{1}B_{2} = B_{3}$, and $x_{r}\tilde{x}_{r} = r \cdot r$, a scalar. Let $\Gamma$ be the set of the numbers $x\tilde{x}$, $x \ne 0$ rational. It is closed under products, $(xy)\widetilde{(xy)} = x(y\tilde{y})\tilde{x} = (x\tilde{x})(y\tilde{y})$, and under inverses, since $x^{-1} = \tilde{x}/(x\tilde{x})$ has $x^{-1}\widetilde{x^{-1}} = 1/(x\tilde{x})$; it contains the squares of the nonzero rationals. Every prime $p$ is in $\Gamma$, by induction on $p$: $2 = (1 + B_{1})(1 - B_{1})$. For odd $p$ put $h = \tfrac12(p - 1)$. The numbers $a^{2}$, $0 \le a \le h$, are pairwise incongruent modulo $p$ (if $p$ divides $(a - a')(a + a')$ with $0 \le a, a' \le h$ then $a = a'$, because $0 < a + a' < p$ otherwise), and so are the numbers $-1 - b^{2}$, $0 \le b \le h$. These are $p + 1$ numbers, so $a^{2} \equiv -1 - b^{2}$ for some $a, b$: $a^{2} + b^{2} + 1 = mp$ with $0 < mp \le 2h^{2} + 1 < p^{2}$, so $1 \le m < p$. Every prime factor of $m$ is smaller than $p$ and lies in $\Gamma$ by induction, so $m \in \Gamma$; $mp$ is the norm of $a + bB_{1} + B_{2}$; so $p = (mp)m^{-1} \in \Gamma$. Hence every positive integer is in $\Gamma$, and every positive rational $a/b = ab/b^{2}$.

\emph{(iii).} By (ii) there is $t \in \Q^{4}$ with $t \cdot t = r \cdot r$; let $s \in \Q^{S}$ have the coordinates of $t$ on $S$. If $s \ne r$, $m = r - s$ has $m \cdot m > 0$ and $m \cdot m = r \cdot r - 2r \cdot s + s \cdot s = -2\,s \cdot m$, so $R(s) = s - 2\frac{s \cdot m}{m \cdot m}m = s + m = r$. $R$ preserves the scalar product and has rational entries, so $W = R(\Q^{S})$ is isometric to $\Q^{S}$, that is, to the standard $\Q^{4}$, by $R$; $R(x) - x$ is a multiple of $m$, so $W \subseteq \Q^{S} + \Q m = \Q^{S} + \Q r$. In the algebra of $\Q^{N}$, $-mxm^{-1} = -(2\,x \cdot m - xm)m^{-1} = x - 2\frac{x \cdot m}{m \cdot m}m$.

\emph{(iv).} Let $x = x_{Tu}$, $y = x_{Tv}$ and $z = y\tilde{x}$. Then $\langle z\rangle_{0} = Tu \cdot Tv = c\,u \cdot v$ (\claimref{synthesis.additive-weights-on-the-rational-projectors-of-two-modes-are-linear}, clause (ii)), and $z\tilde{z} = (x\tilde{x})(y\tilde{y}) = c^{2}(u \cdot u)(v \cdot v)$. With $z = \langle z\rangle_{0} + \beta_{1}B_{1} + \beta_{2}B_{2} + \beta_{3}B_{3}$, $z\tilde{z} = \langle z\rangle_{0}^{2} + \beta_{1}^{2} + \beta_{2}^{2} + \beta_{3}^{2}$. So $(u \cdot u)(v \cdot v) - (u \cdot v)^{2} = \sum_{k}(\beta_{k}/c)^{2}$. Suppose $7$ is a sum of three rational squares. Clearing denominators, $7n^{2} = a_{1}^{2} + a_{2}^{2} + a_{3}^{2}$ with integers $n \ge 1$, $a_{1}, a_{2}, a_{3}$ of greatest common divisor $1$. A square is $0$, $1$ or $4$ modulo $8$, and a sum of three of them is never $7$ modulo $8$ (checked exhaustively). If $n$ is odd, $7n^{2} \equiv 7$ modulo $8$: impossible. If $n$ is even, $a_{1}^{2} + a_{2}^{2} + a_{3}^{2} \equiv 0$ modulo $4$, which forces $a_{1}, a_{2}, a_{3}$ even (squares are $0$ or $1$ modulo $4$), against the greatest common divisor. So $7$ is not such a sum, and the pair of the statement, with $u \cdot u = 1$, $v \cdot v = 7$, $u \cdot v = 0$, lies in no four-dimensional rational subspace similar to the standard one.

\emph{(v).} \emph{Frame functions.} As in the theorem for two modes, clause (i), whose argument does not depend on the size: a symmetric idempotent $E$ is the orthogonal projection onto its range, which has a rational basis of pairwise orthogonal vectors $r_{a}$, and $E = \sum_{a}\Pi_{r_{a}}$; additivity gives $\mu(E) = \sum_{a}\mu(\Pi_{r_{a}})$. Put $F(r) = (r \cdot r)\mu(\Pi_{r})$ for $r \ne 0$ and $F(0) = 0$; $F(\alpha r) = \alpha^{2}F(r)$.

\emph{$F$ is a quadratic form on every standard four-space.} Let $R$ be a rational orthogonal matrix and $S$ a set of four coordinates, $W = R(\Q^{S})$, and $\iota$ the embedding of $\Mat_{4}(\Q)$ as the matrices supported on $S$. Then $\mu'(E) = \mu(R\iota(E)R^{T})$ is defined on the projectors of $\Mat_{4}(\Q)$ and is an additive weight, because $E \mapsto R\iota(E)R^{T}$ maps projectors to projectors and preserves sums and the relation $EE' = 0$. By the theorem for two modes (\claimref{synthesis.additive-weights-on-the-rational-projectors-of-two-modes-are-linear}, clause (iii)), transported by its isomorphism $m$ of $\Cl_{\Q}(2,2)$ with $\Mat_{4}(\Q)$ (clause (i) there), there are $\rho'_{ab} = \rho'_{ba} \in V$ with $\mu'(\Pi_{t}) = \sum\rho'_{ab}t_{a}t_{b}/(t \cdot t)$. Since $R\iota(\Pi_{t})R^{T} = \Pi_{R\iota t}$ and $R\iota t \cdot R\iota t = t \cdot t$, $F(R\iota t) = \sum\rho'_{ab}t_{a}t_{b}$: the restriction of $F$ to $W$ is a quadratic form with values in $V$, $F(w) = \beta_{W}(w,w)$ with $\beta_{W}$ symmetric and bilinear on $W$.

\emph{The candidate.} Put $\rho_{ii} = F(e_{i})$ and $\rho_{ij} = \tfrac12\bigl(F(e_{i} + e_{j}) - F(e_{i}) - F(e_{j})\bigr)$ for $i \ne j$, and $G(r) = F(r) - \sum\rho_{ij}r_{i}r_{j}$. On every standard four-space $G$ is a quadratic form, as a difference of two. $G(e_{i}) = 0$ and $G(e_{i} + e_{j}) = 0$.

\emph{A lemma.} A quadratic form $q(w) = \beta(w,w)$ on a four-dimensional space $W$ that vanishes on three distinct hyperplanes $H_{1}, H_{2}, H_{3}$ is zero. Proof: $\beta$ vanishes on $H_{1} \times H_{1}$ and on $H_{2} \times H_{2}$ by polarisation. Let $K = H_{1} \cap H_{2}$, $x \in H_{1} \setminus K$, $y \in H_{2} \setminus K$, so $W = K \oplus \Q x \oplus \Q y$. For $w = k + \alpha x + \gamma y$: $q(w) = q(k + \alpha x) + q(\gamma y) + 2\gamma\beta(k,y) + 2\alpha\gamma\beta(x,y) = 2\alpha\gamma\beta(x,y)$. $H_{3}$ is not contained in $H_{1} \cup H_{2}$ (a vector space is not the union of two proper subspaces, and $H_{3} \cap H_{i} \ne H_{3}$), so it contains a vector with $\alpha\gamma \ne 0$, and $\beta(x,y) = 0$.

\emph{Induction on the size $k$ of the support of $r$: $G(r) = 0$.} For $k \le 4$, $r$ lies in a coordinate subspace $\Q^{S}$ with $|S| = 4$ (here $N \ge 4$ is used), on which $G$ is a quadratic form (take $R = 1$) whose values at $e_{i}$ and $e_{i} + e_{j}$, $i, j \in S$, vanish; its bilinear form vanishes on the basis, so $G = 0$ on $\Q^{S}$. For $k \ge 5$ choose $S$ inside the support of $r$ and take $W \ni r$ from (iii), $W \subseteq \Q^{S} + \Q r$, so every vector of $W$ has its support inside the support of $r$. For $i$ in the support of $r$ let $\lambda_{i}$ be the $i$-th coordinate restricted to $W$. Each $\lambda_{i}$ is nonzero, because $\lambda_{i}(r) = r_{i} \ne 0$, and the $\lambda_{i}$ have no common zero in $W$ other than $0$, so they span the dual of $W$ and at least four of them are pairwise not proportional. Their kernels are distinct hyperplanes of $W$ whose vectors have support of size at most $k - 1$, so $G$ vanishes on them by the induction hypothesis. By the lemma $G = 0$ on $W$, and $G(r) = 0$.

\emph{Conclusion.} $F(r) = \sum\rho_{ij}r_{i}r_{j}$ for all $r$, which is the first formula; $\mu(E) = \sum_{a}\mu(\Pi_{r_{a}}) = \sum_{i,j}\rho_{ij}\sum_{a}(\Pi_{r_{a}})_{ij} = \sum\rho_{ij}E_{ij}$. The $\rho_{ij}$ are determined by $F$ at $e_{i}$ and $e_{i} + e_{j}$. For $\ell \ge 2$ compose with $m$ of (i), $N = 2^{\ell} \ge 4$: $L(x) = \sum\rho_{ST}m(x)_{ST}$ is linear and restricts to $\mu$; a linear extension $L'$ with $L'(x^{\dagger}) = L'(x)$ is determined by its values on the self-adjoint elements, which are spanned by the projectors $\Pi_{e_{i}}$ and $\Pi_{e_{i} + e_{j}}$, so $L' = L$.

\emph{(vi).} As for two modes. If $\mu \ge 0$ on atoms then $\sum\rho_{ij}r_{i}r_{j} \ge 0$ for every rational $r$; with $c_{k}$ the columns of $m(x)$, $L(xx^{\dagger}) = \sum_{i,j}\rho_{ij}(m(x)m(x)^{T})_{ij} = \sum_{k}\sum_{i,j}\rho_{ij}(c_{k})_{i}(c_{k})_{j} \ge 0$; $L(1) = \mu(1) = 1$ and $L(x^{\dagger}) = L(x)$ by the symmetry of $\rho$. Uniqueness is that of (v). Conversely a state is linear, hence additive on projectors, and $\omega(\Pi) = \omega(\Pi\Pi^{\dagger}) \ge 0$.

\emph{(vii).} By (i) the self-adjoint idempotents of the even subalgebra are the $E_{+} \oplus E_{-}$ with $E_{\pm}$ symmetric idempotents of $\Mat_{N/2}(\Q)$, $N/2 = 2^{\ell-1} \ge 4$. $E_{+} \oplus 0$ and $0 \oplus E_{-}$ are orthogonal projectors of the even subalgebra, so $\mu(E_{+} \oplus E_{-}) = \mu_{+}(E_{+}) + \mu_{-}(E_{-})$ with $\mu_{+}(E) = \mu(E \oplus 0)$, $\mu_{-}(E) = \mu(0 \oplus E)$, each an additive weight on the projectors of $\Mat_{N/2}(\Q)$. By (v) each is linear, and so is $\mu$ on the even subalgebra; uniqueness as in (v), block by block. For $\ell = 2$ the blocks have size two. In size two every ray has exactly one orthogonal ray, so the additive weights on a block are the functions with $f(r) + f(r^{\perp})$ constant, and most of these are not of the form $\sum\rho_{ij}r_{i}r_{j}/(r \cdot r)$; the stored two-valued weight on the even subalgebra of two modes (\claimref{synthesis.no-two-valued-weights-on-the-atoms-of-two-modes-the-even-subalgebra-admits-them}) is one, since a form $\sum\rho_{ij}r_{i}r_{j}/(r \cdot r)$ with values in $\{0, 1\}$ only would be constant on the rays of a block.
